% ============================================================
% PHỤ LỤC C - VIVA FRAMEWORK
% ============================================================

Bộ 3 cẩm nang kèm 90 câu hỏi vấn đáp chi tiết (30 câu/người) được biên soạn và bàn giao tại thư mục \texttt{handbooks/} (gồm \texttt{HANDBOOK\_A.md}, \texttt{HANDBOOK\_B.md}, \texttt{HANDBOOK\_C.md}). Phụ lục này tóm lược các trọng tâm kỹ thuật cốt lõi và ví dụ câu hỏi điển hình của từng thành viên.

\subsection*{Người A -- Agent, Codec, Simulator, Deploy/Gateway}

\textbf{Phải hiểu sâu:} TCP byte stream, framing, partial read, writer queue, SslStream, heartbeat, service probe, reconnect/backoff, retry/dedup ở client, simulator nhiều socket, LAN/Internet topology, least-connections nếu P1.

\begin{vivabox}
Ví dụ câu hỏi cần có trong handbook A:
\begin{enumerate}
    \item Tại sao một lần \texttt{ReadAsync} không thể coi là một message?
    \item Nếu 4 byte length bị chia thành hai lần đọc thì FrameCodec xử lý thế nào?
    \item Tại sao cần giới hạn max frame size?
    \item Heartbeat dùng để giải quyết vấn đề gì mà \texttt{socket.Connected} không đủ?
    \item Backoff giúp tránh reconnect storm ra sao?
    \item Agent Simulator có gì ``thật'' và có gì ``giả''?
    \item Khi Collector chết ở P1, Agent cũ có được migrate socket không? Vì sao blueprint không yêu cầu zero-downtime?
\end{enumerate}
\end{vivabox}

\subsection*{Người B -- Central Server, Database, Rule Engine, Auth Server, Snapshot}

\textbf{Phải hiểu sâu:} accept nhiều client bằng async I/O ThreadPool, session registry, PBKDF2 băm mật khẩu + salt, server-side RBAC, commit ACK, liveness deadline bằng monotonic clock, streak/hysteresis, Incident dual-state, transaction boundary, rowversion OCC, telemetry dedup, retention/restore, internal Collector routing nếu P1.

\begin{vivabox}
Ví dụ câu hỏi cần có trong handbook B:
\begin{enumerate}
    \item Vì sao ACK metrics nên gửi sau một điểm commit xác định?
    \item Hai Admin cùng ACK một Incident tạo race condition như thế nào?
    \item Transaction khác lock trong code ở đâu?
    \item rowversion giúp phát hiện conflict ra sao?
    \item Vì sao cần tách ConditionState khỏi Incident workflow?
    \item Hysteresis khác threshold đơn như thế nào?
    \item Nếu Central restart, dữ liệu nào bắt buộc phải còn và dữ liệu nào có thể mất theo TTL?
\end{enumerate}
\end{vivabox}

\subsection*{Người C -- Admin WinForms, RPC Client, Realtime Cache, Concurrency Test, Collector App}

\textbf{Phải hiểu sâu:} Chu trình WinForms STA, điều phối luồng an toàn (UI marshal bằng \texttt{Invoke}, \texttt{BeginInvoke}, \texttt{InvokeAsync}), đối soát RPC qua \texttt{requestId}, xử lý timeout, đồng bộ cache snapshot/delta, chống race condition khi đa Admin thao tác, cô lập client chậm (\textit{slow-client isolation}) và cấu trúc Collector App (binary P1 dùng chung).

\begin{vivabox}
Ví dụ câu hỏi cần có trong handbook C:
\begin{enumerate}
    \item Authentication và authorization khác nhau thế nào trong LabGuard?
    \item Vì sao ẩn nút trên UI không đủ để bảo vệ quyền RESOLVE?
    \item Snapshot và delta dùng khi nào?
    \item Nếu Admin bị mất một eventSeq rồi reconnect thì phải làm gì?
    \item Vì sao network callback không nên update WinForms control trực tiếp từ background thread?
    \item Slow Admin client có thể ảnh hưởng subscriber khác thế nào và cách cô lập?
    \item Readiness của Collector khác liveness/health check ở điểm nào?
\end{enumerate}
\end{vivabox}

\subsection*{Mẫu trả lời vấn đáp 6 bước}

Khi trả lời một module mình làm, mỗi thành viên nên đi theo thứ tự:
\begin{enumerate}
    \item \textbf{WHAT}: thành phần này là gì?
    \item \textbf{WHY}: nó giải quyết vấn đề gì trong LabGuard?
    \item \textbf{FLOW}: dữ liệu đi từ đâu đến đâu, qua class/API nào?
    \item \textbf{FAILURE}: nếu mạng lỗi/concurrent/invalid input thì sao?
    \item \textbf{VERIFY}: test/log nào chứng minh code hoạt động?
    \item \textbf{TRADE-OFF}: vì sao chọn cách này thay vì cách khác đơn giản/phức tạp hơn?
\end{enumerate}
